\section{Method: A Default Pipeline with Stated Inputs}\label{sec:method}

The pipeline turns an image and a task into a final set, a task answer, a status and an alarm score. Every input it uses comes from one of four named sources: \src{T} the task specification given by the user (classes, alias map, matching rule, loss); \src{S} the shipped configuration, a fixed file released with the code (\texttt{config/default\_pipeline.json}) whose values were chosen on our evaluation data and are carried unchanged to new images; \src{O} quantities computed from the current outputs; and \src{L} optional labels. No step requires the user to know $p_1$, a partner's quality or the current F1. Algorithm~\ref{alg:pipeline} gives the steps in order. Steps 1 to 6 repair the report, the optional label rule of Step 7 judges between configurations, and the alarm returned in Step 8 detects errors. Section~\ref{sec:method:clicks} adds a second tier for clicks that combines several models, and Section~\ref{sec:method:revise} describes how the model behind these steps was revised.

\begin{algorithm}[t]
\caption{Default reporting pipeline}\label{alg:pipeline}
\small
\begin{algorithmic}[1]
\Require image $x$; task $(\mathcal{C}_\tau,\alpha_\tau,\rho_\tau,\Omega_\tau,\ell_\tau)$ \src{T}; budget $B$ per item \src{T}; shipped configuration $\Theta$ \src{S}; optional labels \src{L}
\State \textbf{Dispatch.} From $(\tau,B,\Theta)$ select the reporting interface, setting and call schedule, and send the selected request: counts, the list-first request (record list first, then counts ``for the records you listed''; allowed class names; row-by-row scan; native coordinate frame; no overlays); sets, the list or regions interface that $\Theta$ gives the model and loss, $K$ times with $K$ set by the loss; clicks, one model's two-stage zoom, or Algorithm~\ref{alg:clicks} if $B$ covers its worst case; keep every attempt and its cost \src{T, S}
\State \textbf{Decode.} For each response, $H\mapsto(Z^{\mathrm{all}},Z,d,\sigma)$: audit records, delivered records, stated count, status L1 to L3 with the four degeneracy flags \src{O}
\State \textbf{Recover, at most once per call.} Transport failure (L1): resend the same request. Flagged response (L2 with no recoverable record, or a degeneracy flag): one follow-up that attaches the previous answer; keep it if unflagged, else keep the first response and its flags \src{S, O}
\State \textbf{Views.} Clicks: two-stage zoom on native pixels (ReGround, \citealp{li2025screenspotpro}). Sets and counts: the whole image, or the views of an external view planner if one is configured, one call per view, records mapped to image coordinates \src{S, O}
\State \textbf{Construct.} Merge the reports of all views; map the delivered records of all calls to the final set $S$ with the construction $\Theta$ assigns to the loss (Table~\ref{tab:defaults}); a view confirms or rejects only records it contains \src{S}
\State \textbf{Read.} $\hat y\gets$ list count of $S$; if the list behind $S$ is flagged and $d\neq\varnothing$, $\hat y\gets d$ with status \texttt{count\_from\_statement}; clicks: the final point \src{O}
\State \textbf{Optional labels.} $\ge3$ labelled images: switch to a menu configuration only if its mean paired advantage minus one standard error is positive; 5 to 10 images: may fit the record gate \src{L}
\State \Return $S$, $\hat y$, readout branch, status $\sigma$, alarm score (disagreement with a second report, from a second model if configured), total cost of all attempts
\end{algorithmic}
\end{algorithm}

\begin{table*}[tp]
\centering\small\setlength{\tabcolsep}{4pt}
\caption{Shipped defaults for one model. \emph{Transferred construction}: the best construction on the other five sources, applied to the evaluated source (Table~\ref{tab:transfer}). \emph{Final default}: what the pipeline ships. They coincide for set error on over-reporting families and differ elsewhere; the last column gives the reason.}\label{tab:defaults}
\begin{tabular}{@{}>{\raggedright\arraybackslash}p{2.7cm}>{\raggedright\arraybackslash}p{3.1cm}>{\raggedright\arraybackslash}p{2.7cm}>{\raggedright\arraybackslash}p{6.5cm}@{}}\toprule
Loss and family & Transferred construction & Final default & Evidence and reason \\\midrule
$E$, over-reporting families (detection, layout, screen parsing) & three-call unanimity, averaged coordinates (chosen in all six folds) & the same & Significantly better in 9 of 15 cells, significantly worse in none, not significant in 6. \\
$E$, under-reporting family (dense points, FSC-147) & three-call unanimity, averaged coordinates & single call, no deletion & The transferred construction is worse in 2 of 3 cells. Deleting unconfirmed records raises $E$ when their marginal matching rate $p_1$ is above 1/2 (Proposition~\ref{thm:intersect}); on FSC-147, $p_1$ is 0.39 to 0.90 and deletion hurts in 7 of 9 pairs (Table~\ref{tab:intersect}). \\
F1 & three-call union (five folds) or majority (one fold) & single call; on detection sources, interface and setting chosen per model (GPT-6.1: regions, lowest; Sonnet 5.5: regions, default; Gemini 3.8: list, default) & The transferred construction is better in 2 of 18 cells, worse in 4 and not significant in 12. A configuration chosen on three detection sources is better on the fourth in 5 of 12 cells and worse in 1 (Table~\ref{tab:config}). \\
Count & not part of the transfer & single call; list count, with the stated count as declared fallback & Table~\ref{tab:readout}. \\\bottomrule
\end{tabular}
\end{table*}

\paragraph{Inputs \src{T, S}.} The task fixes the class vocabulary $\mathcal{C}_\tau$, the alias map $\alpha_\tau$ (for example \{diamond, rhombus, rotated square\} $\mapsto$ square), the matching rule $\rho_\tau$ of Table~\ref{tab:contract}, the ignore region $\Omega_\tau$ and the loss $\ell_\tau$. The loss selects the construction in Step 5. The shipped configuration fixes everything else that the algorithm calls a value: the prompt, the reasoning setting per model, the four flag thresholds, the recovery request, the view settings, the construction per loss, the alarm partner and the alarm thresholds. A model known to write degenerate lists at its lowest setting (Gemini 3.8) is run at its default setting, which at a fixed 16k output cap lowers its count readout error from 52.8 to 24.5 (Table~\ref{tab:setting}).

\paragraph{Steps 1 and 2: dispatch and decode.} Step 1 is a dispatcher, not one fixed request: what is shared across tasks is the chain of records, losses and readouts, not one output format. For set tasks scored by F1 the shipped configuration may select the regions interface, in which the model returns one row per region of a fixed coordinate registry (Appendix~\ref{app:protocol}). The count request asks for one JSON object in which the list of records comes first and \texttt{counts\_by\_type} follows, declared ``for the shapes you listed'' (Appendix~\ref{app:protocol} gives the prompt). The allowed class names are listed in the format; the model scans row by row from top to bottom and uses its native coordinate frame (0 to 1000 for GPT and Gemini; pixels for Sonnet, with the image pre-resized so that the provider does not resize it again). No grid, ticks or markers are drawn. The decoder keeps every recoverable record in the audit collection $Z^{\mathrm{all}}$ and passes the valid ones, after the alias map, to the delivered collection $Z$ (Section~\ref{sec:setup}). It sets four degeneracy flags on the list: more than 10\% of records without a valid class; a list length outside $[0.5,1.5]$ times the stated total; more than 10\% repeated coordinates; and the finish reason ``length''.

\paragraph{Step 3: bounded recovery.} The two kinds of failure get different recoveries. A transport failure, in which the call did not return, is a property of the connection and is resent unchanged. A flagged response is a property of how the model wrote this report, and such errors tend to recur, consistent with Section~\ref{sec:theory:behaviour}: in our detection experiments, resending a request whose flag had already repeated cleared the flag in 14.7\% of cases, while changing the request cleared it in 58.8\%. The change was a follow-up that attached the previous answer and asked for the complete answer under the same contract, or, for reports organised by image region, a new request for the flagged regions only. The pipeline sends the first kind of follow-up, once. If the follow-up still carries a flag, the first response is kept with its flags, and Step 6 decides how to read it. Every attempt enters the cost. The replay at the end of this section reports how often recovery was triggered and what the stored follow-ups changed.

\paragraph{Step 4: views.} For clicks we use a re-grounding strategy in the style of the ReGround method that \citet{li2025screenspotpro} evaluate. The first call gives a point on the whole screenshot. A crop one third of the screenshot's width and height, centred on that point and shifted to stay inside the image, is sent at native pixels, and the second call gives the point in the crop. If the second call gives no point, the first point is kept; if the first call gives no point, there is nothing to zoom on and the answer is no point, which counts as a miss. Using agreement to decide when to zoom costs more and is not more accurate (Section~\ref{sec:res:view}). For sets and counts, the pipeline does not plan views itself: it works on whatever reports it is given. By default it sends the whole image. When a \emph{view planner} is configured, for example a tiling scheme or the companion paper's safe decomposition, which places every target wholly inside some view, with neighbouring views overlapping by one object and no view having fewer tokens per object than the whole image, and cannot lower expected recall under that paper's monotonicity assumptions \citep{shi2026zoomin}, each view is one call, its records are mapped to image coordinates, and Step 5 merges them. Leaving the planner and its trigger to the configuration is a measured choice. We computed the apparent size of every annotated object in the whole-image calls on SCP-dense, DocLayNet, DOTA and SKU-110K, using the input geometry measured in the companion paper and its released calibration data (32-pixel tokens for GPT-6.1, 28-pixel tokens for Sonnet 5.5, a fixed grid of about 1,090 tokens for Gemini 3.8). The median object spans 0.62 to 1.43 tokens across its shorter side (the companion paper measures size as the square root of the area, which is never smaller). Below that paper's design value of $S^\star=1.5$ fall 55\% to 79\% of objects on the three detection sources and every shape on SCP-dense; yet GPT-6.1 and Sonnet 5.5 write nearly exact lists on SCP-dense. The companion paper likewise finds that these model families reach half recall at well under one token per object side. A fixed threshold would decompose images whose whole-image report is already complete, so when to call a planner is a property of the model and the task that the configuration supplies. These sizes come from the annotations and inform only this design choice; the pipeline reads no annotation at run time. Section~\ref{sec:res:view} reports how the pipeline handles reports from four crops, and how our own tiling arm, which did not meet the safe condition, fared on large objects.

\paragraph{Step 5: construction.} The construction is chosen per loss (Table~\ref{tab:defaults}). Records of different calls are associated without labels (Table~\ref{tab:contract}): boxes of the same class with IoU $\ge0.5$, screen elements at IoU $\ge0.5$ without class, and points of the same class within the task's tolerance. On over-reporting task families and set error, three calls are made and a record is kept only if all three report it, with its coordinates averaged over the three. Everywhere else a single call is delivered without deletion. Over- or under-reporting is a property of the task family, shipped with the configuration from our evaluation; the user does not supply it. Coordinate averaging is on by default because its contribution is positive in 37 of 54 base/partner pairs and negative in 13 (Table~\ref{tab:intersect}); the deletion-only variant is in the menu. Reports of several views are first merged into one report: the union of their records, where a record is dropped if it is associated with a record already kept from another view (a duplicate in an overlap). When reports from different views confirm one another, a view can confirm or reject only records that lie wholly inside it; a record that no partner view contains is kept, because that view carries no evidence about it (Section~\ref{sec:theory:behaviour}). A call that still has no readable answer after recovery counts as an empty report, so three-call unanimity keeps nothing on that image; this is the construction we evaluated, and dropping such a call from the support count is in the menu.


\paragraph{Step 6: readout.} The answer is read from the final set. For counts it is the list count of $S$, the number of records of each class. When the list behind $S$ carries a degeneracy flag and the response stated counts, the answer falls back to the stated count $d$, and the status says so (\texttt{count\_from\_statement}); a flagged list without stated counts keeps the list count and its flag. With several views, each response states counts for its own crop, and overlapping views would count shared objects twice; stated counts are therefore never combined across views, and a flagged merged list keeps the list count, with the status saying that no stated count was used. The fallback changes only the count: $S$ itself is not repaired, so the count and the set each carry their own status. Proposition~\ref{prop:readout} explains why the rule is a fallback and not a selector: a large gap between list count and stated count shows that the two readouts of one response are inconsistent, but by itself it does not say which one is closer to the truth. The flags are the evidence that the list is broken.

\paragraph{A second model \src{S}, optional.} When a second model is configured, its default role is the alarm of Step 8: its answer is compared with the user's, and items where they disagree by more than a threshold are flagged. The partner is fixed by a rule (GPT-6.1 and Sonnet 5.5: Gemini 3.8; Gemini 3.8: Sonnet 5.5), which picks the cheapest other model per click call (Appendix~\ref{app:protocol}), and the threshold $\delta$ is a point the user picks on the risk-coverage curve (Figure~\ref{fig:alarm}). Using the second model to confirm records is a menu option. The cross-model construction transferred from the other sources improves set error in 12 of 18 transfer cells and F1 in 10 of 18, and is worse in 2 and 3 cells; all of the worse cells are cells in which the user's own model is the strongest single model for that source and loss (Table~\ref{tab:transfer}). Because the user may not know whether that is the case, cross-model confirmation is chosen with labels or a shipped family-level ranking, not as a zero-label default.

\paragraph{Step 7: optional labels \src{L}.} Labels are never required. With fewer than three labelled images the pipeline keeps the shipped default and reports the labelled error. With three or more, it switches to another configuration from the menu only if that configuration's mean advantage over the labelled images, paired by image, minus one standard error is positive. For the decision to delete unconfirmed records, this rule predicts the better decision on the other images in 86\% to 88\% of draws with 3 to 10 labelled images, against 85\% for the family default alone (Section~\ref{sec:res:set}). With 5 to 10 labelled images it may also fit the record gate, a logistic selector on output features (Section~\ref{sec:res:select}). Certifying the accuracy of a pipeline on a population is a separate service that needs its own labelled sample.

\paragraph{Step 8: delivery.} The pipeline returns the final set $S$, the answer $\hat y$, the readout branch, the status $\sigma$ (L1 to L3, flags, recovery used, views used), the alarm score when a second report exists, and the cost of every attempt.

\paragraph{One entry point, replayed on stored calls.} The repository implements Steps 1 to 6 and 8 of Algorithm~\ref{alg:pipeline}, including merging views and cross-view confirmation, as a single entry point (\texttt{fsr\_pipeline/}) whose branches read only the task specification, the shipped configuration and the responses a client returns; the label rule of Step 7 is a function in the same module, and the record gate is fitted by the analysis code of Section~\ref{sec:res:select}. A replay client serves the stored responses of this paper instead of calling a model. A request for which nothing was stored, such as a follow-up that was never sent, returns ``unavailable'', and the pipeline continues on its documented branch. Run on raw responses, the entry point reproduces the component analyses: per-type count error of 1.60 for GPT-6.1, 0.42 for Sonnet 5.5 and 24.52 for Gemini 3.8 at its default setting (52.80 at its lowest; Table~\ref{tab:setting}); click hit rates of 90.0\%, 87.0\% and 87.5\% after two-stage zoom on the 200 ScreenSpot-Pro items; and, with recovery switched off, the set error of the shipped construction in all 18 transfer cells to within 0.07 per image, a difference that comes from the 0.1-pixel coordinate rounding of the component analysis. Follow-ups were stored for a sample of flagged detection calls. In the 18 set cells, 171 image runs had at least one flagged call; a stored follow-up cleared the flag in 24 of them, and replaying these changes set error by at most 0.55 per image in any cell. No follow-ups were stored for the count experiments, where 51 runs were flagged, so their results are those of the first call. Synthetic boundary cases go through the same code, and each produces an answer, a readout branch and a status: an empty or unreadable response, truncated JSON, a list without stated counts, a degenerate list with and without stated counts, two transport failures in a row, one unreadable call among three, and a click with no first point or no point in the crop. The replay output is released with the analysis outputs (\texttt{pipeline\_replay.json}).

\subsection{Clicks from several models}\label{sec:method:clicks}

Algorithm~\ref{alg:pipeline} delivers one model's two-stage zoom point. When the per-item budget covers the worst case of a second tier, the pipeline offers one built on the three sources of a miss (Proposition~\ref{prop:three}); its mean cost is about 3.2 times that of one zoom. Its run-time inputs are the image, the instruction and the budget; it reads no target label and does not abstain: it returns no point only when no call returned one. The models, their settings, the distance $\delta$, the judge and the number of compared clusters come from the shipped configuration \src{S}.

\begin{algorithm}[t]
\caption{Clicks from several models}\label{alg:clicks}
\small
\begin{algorithmic}[1]
\Require screenshot $x$, instruction and budget $B$ \src{T}; models, settings, per-call output caps, $\delta$, judge, source order, $m_{\max}$ \src{S}
\State \textbf{Tier.} Run this tier only if $B$ covers its worst case at the per-call caps (15 calls besides technical retries: three first clicks, three zooms, three scores, at most six comparisons); otherwise run one model's two-stage zoom \src{T, S}
\State \textbf{Generate.} GPT-6.1 and Sonnet 5.5 at their default setting and Gemini 3.8 at its lowest setting each give a two-stage zoom point (Step 4 of Algorithm~\ref{alg:pipeline}; a failed zoom keeps the first click, a failed first click gives no candidate); each point is registered as a candidate with an identifier bound to its coordinates on $x$, its model, setting and parent call \src{S, O}
\State \textbf{Detect.} Clusters are the connected components of the graph that joins candidates within $\delta=1.5\%$ of the longer side; two or more clusters raise the alarm \src{O}
\State \textbf{Represent.} The judge (Gemini 3.8, lowest setting) sees a native crop with each candidate marked and returns a hit score $p\in[0,1]$ (a missing or unreadable score counts as $-1$); each cluster's representative is its member with the highest $p$, ties broken by the source order GPT-6.1, Sonnet 5.5, Gemini 3.8, fixed before any comparison \src{S, O}
\State \textbf{Compare (alarm only).} Rank clusters by (representative's $p$, cluster size, earliest source in the cluster); for each pair among the first $m_{\max}=3$, the judge sees both representatives marked with their identifiers on the downscaled screenshot and on a native crop around each, and names the one that hits the target; each pair is asked in both orders, and a representative wins only if both answers name it; a failed, undecided or split pair gives no win \src{S, O}
\State \Return the representative of the cluster with most wins, with ties among those clusters broken by the ranking of Step 5; if no cluster has a win, the first cluster of that ranking; if no call gave a point, no answer (a miss); the alarm and the cost of all attempts
\end{algorithmic}
\end{algorithm}

Algorithm~\ref{alg:clicks} gives the steps, organised by Proposition~\ref{prop:three}. Generation enlarges the candidate menu, which is the only way to lower $F$; Section~\ref{sec:res:choose} measures the rescue rate of each source on the others' misses (Proposition~\ref{prop:coverage}). The representative step fixes which member of each cluster can be output, and the comparison chooses among representatives, so these two steps divide the remaining misses between $L$ and $W$. The comparison runs only on items with an alarm, where the measured miss rate of the judge's score is highest (Proposition~\ref{prop:switch}); items whose candidates form one cluster keep that cluster's representative. The judge's score is used only to rank; we do not treat it as a calibrated probability. The registry makes the point that is scored, drawn, compared and returned the same point: every step refers to a candidate by its identifier, and the code asserts that the coordinates drawn for the judge equal those returned. If a call fails, the item proceeds with the remaining candidates, and it counts as a miss only when no call returned a point. The budget rule uses the worst case, at per-call output caps of 16,384 tokens at the lowest setting and 32,768 at the default, because a mean cost is not a per-item guarantee; technical retries can add to either. On 600 new screenshots the tier costs \$0.0273 per item on average (9.45 calls) against \$0.0085 (2 calls) for GPT-6.1's two-stage zoom (Section~\ref{sec:res:choose}), and no budget was imposed in our runs. The rules are implemented as one deterministic function over stored or live responses (\texttt{gui/code/decision.py} in the release), with request builders for each step; with the evaluation script, run on the stored responses, it regenerates Table~\ref{tab:choose}.

\subsection{Revising the model from failed predictions}\label{sec:method:revise}

The theory was built by iteration, and the iteration is part of the method. A model is stated, a consequence is derived and tested, and when a prediction fails we record three things: which assumption failed and how we know; how the revised assumption was checked directly; and what new consequence the revision predicts, and whether it holds. Table~\ref{tab:revisions} lists the five revisions that shaped the final theory. Each row names the kind of evidence: an exact identity, a controlled comparison of interventions on the same images, or an association observed across cells or records.

\begin{table*}[tp]
\centering\small\setlength{\tabcolsep}{4pt}
\caption{Revisions of the model. Each row records the failed assumption and the evidence against it, the revised assumption and its direct check, and the new consequence with its test. The kind of evidence is named in each cell.}\label{tab:revisions}
\begin{tabular}{@{}>{\raggedright\arraybackslash}p{3.1cm}>{\raggedright\arraybackslash}p{3.4cm}>{\raggedright\arraybackslash}p{4.0cm}>{\raggedright\arraybackslash}p{4.5cm}@{}}\toprule
Failed assumption & How we know & Revised assumption and direct check & New consequence and test \\\midrule
The intersection threshold, applied to intersections with averaged coordinates, predicts whether confirmation helps. & Wrong sign in 2 of 54 base/partner pairs. & The threshold governs deletion only; coordinate averaging is a separate term (Proposition~\ref{thm:intersect}, Proposition~\ref{prop:geometry}). Exact identity: both terms add up to the observed change on every image of all 54 pairs. & The averaging term can take either sign and can reverse the result. Observed: positive in 37 pairs, negative in 13; it reverses the deletion-only result in two pairs, which accounts for both former exceptions (Table~\ref{tab:intersect}). The deletion term, in turn, can be predicted from $p_1$ on 5 labelled images in 83\% of draws (Section~\ref{sec:res:set}). \\
Similar quality is enough for a second model to help as a confirming partner. & On ScreenParse the three models have similar set error (16.2, 13.3, 16.0), yet among constructions with averaged coordinates the best cross-model construction is not significantly different from the best same-model one ($+1.8\%$, $[-3.0,+7.6]$). This refutes sufficiency only. & A partner helps if it rejects more of the base model's false records than it loses of its true ones (Proposition~\ref{prop:partner}). Direct check across pairs: cross-model partners separate the two better than same-model partners in 32 of 36 comparisons; the 4 exceptions are partners that confirm few of the base model's true records. & Conditional consequence: a base with many true and few false records gains only if the partner keeps almost all of its true records ($T(1-h_+^B)$ dominates $G$), and a weaker partner usually keeps fewer, so cross-model confirmation is expected to do worst where the base model is strongest. Observed association across transfer cells: all 5 significantly worse cells are cells where the base model is strongest, and 20 of the 24 cells where it is not strongest improve (Table~\ref{tab:transfer}). \\
A count written first, or given in the prompt, pulls the list length like a quota for every model. & Giving a count moves list length with slope at most 0.004 for GPT-6.1 and Sonnet 5.5 on SCP and SCP-dense (controlled comparison). & The pull depends on the model and on how reliable its enumeration is. Direct check: slope 0.17 to 0.20 on PixMo, and 1.63 for Gemini 3.8 on SCP-dense, where its lists degenerate. & Asking for the list before the counts should not hurt and should help where enumeration is reliable. Observed in a controlled comparison under the same mapped readout: list-first lists beat count-first lists for all three models ($-39\%$, $-92\%$, $-25\%$; Table~\ref{tab:readout}). The list-first prompt also adds the row-by-row scan and the allowed class names, so this is a test of the protocol, not of order alone. \\
Gemini 3.8 has a coordinate spread that a different coordinate convention reduces. & Pixel versus grid convention: $+2\%$, n.s.\ (controlled comparison). & Current models make no coordinate-frame error on SCP. Direct check: systematic offset at most 0.3 pixels, slope 1.000, spread 0.9 to 1.5 pixels, against a 30-pixel tolerance. & References that correct coordinates (grids, ticks) should not improve localization. Observed in controlled comparisons: localization changes by at most 0.23 pixels in every arm. Any effect of an overlay must come through how the model writes its list (next row). \\
A grid helps Gemini 3.8 by improving localization. & The grid lowers set error by 3.5 per image at the lowest setting, but localization does not change, and at the default setting the grid makes detection slightly worse (controlled comparisons). & The grid acts on the record format: its gain falls on records that omit their class field. Association within lists: a record omits its class more often when the previous record has the same class (logistic regression controlling for position in the list: $+0.70$ $[0.57,0.84]$ in log-odds, arm fixed effects, intervals clustered by image). This path explains only a small part of the grid's gain. & An order that interleaves classes should lower omissions, and grouping by class should raise them. Observed in a controlled comparison: a row-by-row scan lowers set error by 11.5 per image against 3.5 for the grid, and grouping by class raises it (Section~\ref{sec:res:degen}). \\\bottomrule
\end{tabular}
\end{table*}
